
The following examples are the saddle point problems of rational polynomials studied by \textcite{zhou.wang.ea:saddle}, which we interpret here as two-player zero-sum games where each player controls multiple variables.

\begin{example}\label{ex.zhou19.5.3.i}
A two-player zero-sum game in $2 \times 2$ variables on the unit hypercube
\[
x_i \in [0,1]^2 \quad \forall i \in \{1,2\}.
\]
The rational utility function \cite[Example 5.3 (i)]{zhou.wang.ea:saddle} is given by
\[
u_1(x_1, x_2) =
\frac{
- x_{11}^2 - x_{12}^2 + x_{21}^2 + x_{22}^2
}{
x_{11} + x_{22} + 1
}
\]
This game has a pure Nash equilibrium at
\begin{align*}
x_1^\star = (0, 0),\, x_2^\star = (0, 0).
\end{align*}
\end{example}


\begin{example}\label{ex.zhou19.5.3.ii}
A two-player zero-sum game in $3 \times 3$ variables on the unit hypercube
\[
x_i \in [0,1]^3 \quad \forall i \in \{1,2\}.
\]
The rational utility function \cite[Example 5.3 (ii)]{zhou.wang.ea:saddle} is given by
\[
u_1(x_1, x_2) =
\frac{
- \sum_{i=1}^{3} (x_{1i} + x_{2i}) - \sum_{\substack{i=1 \\ i < j}}^{3} \sum_{j=1}^{3} (x_{1i}^2 x_{2j}^2 - x_{2i}^2 x_{1j}^2)
}{
x_{11}^2 + x_{22}^2 + x_{13} x_{23} + 1
}
\]
This game has no pure Nash equilibria.
\end{example}



